% Part: many-valued-logic
% Chapter: syntax-and-semantics
% Section: connectives

\documentclass[../../../include/open-logic-section]{subfiles}

\begin{document}

\olfileid{mvl}{syn}{con}

\olsection{ژبې او نښلوونکي}

کلاسيک بياني منطق، او ډېر نور منطقونه، د \emph{بياني ثابتونو} او
\emph{نښلوونکو} يو سټ کاروي۔ د بېلګې په توګه، لاندې نښې د بنسټيزو
نښو په توګه کاروو:
\begin{enumerate}
\tagitem{prvFalse}{د !!{falsity} بياني ثابت~$\lfalse$۔}{}
\tagitem{prvTrue}{د !!{truth} بياني ثابت~$\ltrue$۔}{}
\item منطقي نښلوونکي:
  \startycommalist
  \iftag{prvNot}{\ycomma $\lnot$ (نفي)}{}%
  \iftag{prvAnd}{\ycomma $\land$ (عطف)}{}%
  \iftag{prvOr}{\ycomma $\lor$ (فصل)}{}%
  \iftag{prvIf}{\ycomma $\lif$ (!!{conditional})}{}%
  \iftag{prvIff}{\ycomma $\liff$ (!!{biconditional})}{}%
\end{enumerate}
\iftag{defNot,defOr,defAnd,defIf,defIff,defTrue,defFalse,defEx,defAll}{%
د پورته بنسټيزو نښلوونکو تر څنګ، هغه نښې هم کاروو چې د لنډيزونو
په توګه تعريف شوې دي، لکه
\startycommalist
  \iftag{defNot}{\ycomma $\lnot$ (نفي)}{}%
  \iftag{defAnd}{\ycomma $\land$ (عطف)}{}%
  \iftag{defOr}{\ycomma $\lor$ (فصل)}{}%
  \iftag{defIf}{\ycomma $\lif$ (!!{conditional})}{}%
  \iftag{defIff}{\ycomma $\liff$ (!!{biconditional})}{}%
  \iftag{defFalse}{\ycomma $\lfalse$ (!!{falsity})}{}%
  \iftag{defTrue}{\ycomma $\ltrue$ (!!{truth})}.}{}

په څو ارزښته منطقونو کښې هم همدغه نښلوونکي کارېږي۔ خو ډېر ځله
ګټوره وي چې د بېلګې په توګه د عطف بېلابېلې بڼې په يوه منطق کښې
شاملې کړو؛ بيا د دغو بڼو د بېلولو دپاره بېلو نښو ته اړتيا وي۔
ځينې څو ارزښته منطقونه داسې نښلوونکي هم لري چې په کلاسيک منطق
کښې ئې سيال نشته۔ له همدې امله دلته له معمول څخه لږ عمومي
تعريف کاروو۔

\begin{defn}
  يوه \emph{بياني ژبه} د \emph{نښلوونکو} له يوه سټ~$\Lang L$ څخه
  جوړه وي۔ هر نښلوونکے~$\star$ يو ټاکلے شمېر ځايونه لري؛ هغه
  نښلوونکے چې~$n$ ځايونه لري، \emph{$n$-ځايه} بلل کېږي۔
  د~$0$ ځايونو نښلوونکو ته \emph{ثابتونه} هم ويل کېږي؛
  د~$1$ ځاے نښلوونکي \emph{يو ځايه}، او د~$2$ ځايونو
  نښلوونکي \emph{دوه ځايه} بلل کېږي۔
\end{defn}

\begin{ex}
  د بياني منطق معياري ژبه~$\Lang L_0$ له لاندې نښلوونکو
  (او د هر يوه له ټاکلي شمېر ځايونو) څخه جوړه ده:
  $\lfalse$~($0$)،
  $\lnot$~($1$)،
  $\land$~($2$)،
  $\lor$~($2$)،
  $\lif$~($2$)۔ ډېری هغه منطقونه چې موږ ئې څېړو، همدا ژبه
  کاروي۔ ځينې منطقونه د دود له مخې د ځينو نښلوونکو دپاره
  بېلې نښې کاروي۔ د بېلګې په توګه، د حاصل‌ضرب په منطق کښې
  د عطف نښه ډېر ځله د~$\land$ پر ځاے~$\odot$ وي۔ کله ناکله
  د نوي عملګر زياتول هم ګټور وي، لکه د ټاکلتيا عملګر~$\triangle$
  ($1$-ځايه)۔
\end{ex}

\end{document}
